% ============================================================================
% kapitel/frozen.tex
% ============================================================================

\section{Die eingefrorenen Funktionen}

Die Simulation hat zwei eingefrorene Funktionen.

\subsection{Fluktuationen (Vakuum)}

Die Fluktuationen wirken nur im Vakuum.

Sie sind die Quelle der Strukturbildung.

\subsection{Energiesenke (Strukturen)}

Die Energiesenke wirkt nur in Strukturen.

Sie entfernt Energie aus den Strukturen.

\subsection{Die Balance}

Die Balance zwischen Fluktuation und Energiesenke ist gefunden.

\[
\text{Fluktuation} = \text{Energiesenke}
\]

Das System ist stabil.

\subsection{Die Bedeutung}

Die eingefrorenen Funktionen sind das Herz der Simulation.

Sie zeigen, dass die IWT funktioniert.

Sie sind der Beweis für die Theorie.